\section{Detailed Mathematical Proofs and Theoretical Analysis}
\label{sec:appendix_proofs}

In this section of the appendix, we present the complete, self-contained mathematical proofs and theoretical analyses for the theorems and concepts stated in Section 3 of the main body.

\subsection{Proof of Theorem~\ref{thm:rademacher} (Empirical Rademacher Complexity of Fourier Trajectories)}
\label{proof:thm_rademacher}

\begin{theorem}[Empirical Rademacher Complexity of Fourier Trajectories]
Let $\mathcal{H}_F$ be the class of Fourier ensembling trajectories of spectral cutoff $F$ with $\ell_1$-bounded coefficients $\|\theta\|_1 \le C_0$. The empirical Rademacher complexity of $\mathcal{H}_F$ over the network depth coordinate set $Z = \{z_1, \dots, z_L\}$ of size $L$ satisfies:
\begin{equation}
    \widehat{\mathcal{R}}_L(\mathcal{H}_F) \le C_0 \sqrt{\frac{2 \ln(4F+2)}{L}}
\end{equation}
\end{theorem}

\begin{proof}
Any trajectory function $h \in \mathcal{H}_F$ can be represented as an inner product $h(z) = \langle \theta, \Psi(z) \rangle$, where:
\begin{equation}
    \Psi(z) = [1, \cos(2\pi z), \sin(2\pi z), \dots, \cos(2\pi F z), \sin(2\pi F z)]^T \in \mathbb{R}^{2F+1}
\end{equation}
By definition, the empirical Rademacher complexity of $\mathcal{H}_F$ over the coordinate sample $Z$ of size $L$ is:
\begin{equation}
    \widehat{\mathcal{R}}_L(\mathcal{H}_F) = \mathbb{E}_{\sigma} \left[ \sup_{h \in \mathcal{H}_F} \frac{1}{L} \sum_{l=1}^L \sigma_l h(z_l) \right]
\end{equation}
where $\sigma_l \in \{-1, 1\}$ are independent Rademacher random variables. Substituting $h(z_l) = \langle \theta, \Psi(z_l) \rangle$, we obtain:
\begin{equation}
    \widehat{\mathcal{R}}_L(\mathcal{H}_F) = \mathbb{E}_{\sigma} \left[ \sup_{\|\theta\|_1 \le C_0} \left\langle \theta, \frac{1}{L} \sum_{l=1}^L \sigma_l \Psi(z_l) \right\rangle \right]
\end{equation}
By H{\"o}lder's inequality, for any vectors $u, v$, we have $\langle u, v \rangle \le \|u\|_1 \|v\|_\infty$. Applying this to the inner product yields:
\begin{align}
    \widehat{\mathcal{R}}_L(\mathcal{H}_F) &\le \mathbb{E}_{\sigma} \left[ \sup_{\|\theta\|_1 \le C_0} \|\theta\|_1 \left\| \frac{1}{L} \sum_{l=1}^L \sigma_l \Psi(z_l) \right\|_\infty \right] \\
    &\le C_0 \mathbb{E}_{\sigma} \left[ \max_{j \in \{1, \dots, 2F+1\}} \left| \frac{1}{L} \sum_{l=1}^L \sigma_l \Psi_j(z_l) \right| \right]
\end{align}
where $\Psi_j(z)$ represents the $j$-th component function of $\Psi(z)$. Since the range of each component function is bounded, i.e., $\Psi_j(z) \in [-1, 1]$ for all $z \in [0, 1]$, each coordinate sum $X_j = \sum_{l=1}^L \sigma_l \Psi_j(z_l)$ is a sum of independent, zero-mean, sub-Gaussian random variables with parameter at most $1$.

Consequently, the term $\frac{1}{L} \sum_{l=1}^L \sigma_l \Psi_j(z_l)$ is sub-Gaussian with parameter $1/L$. Taking the maximum over both the positive and negative coordinates to evaluate the absolute value, we are maximizing over $2(2F+1) = 4F+2$ sub-Gaussian variables. Applying Massart's Finite Lemma to this finite set of size $4F+2$ yields:
\begin{equation}
    \mathbb{E}_{\sigma} \left[ \max_{j \in \{1, \dots, 2F+1\}} \left| \frac{1}{L} \sum_{l=1}^L \sigma_l \Psi_j(z_l) \right| \right] \le \sqrt{\frac{2 \ln(4F+2)}{L}}
\end{equation}
Substituting this back into our inequality gives the exact complexity bound:
\begin{equation}
    \widehat{\mathcal{R}}_L(\mathcal{H}_F) \le C_0 \sqrt{\frac{2 \ln(4F+2)}{L}}
\end{equation}
This completes the proof.
\end{proof}


\subsection{Proof of Theorem~\ref{thm:rademacher_joint} (Joint Multi-Task Trajectory Complexity)}
\label{proof:thm_rademacher_joint}

\begin{theorem}[Joint Multi-Task Trajectory Complexity]
The joint empirical Rademacher complexity of the multi-task trajectory class $\mathcal{H}_F^{\text{joint}}$ over network depth $L$ is strictly independent of the task count $K$ inside the logarithm, satisfying:
\begin{equation}
    \widehat{\mathcal{R}}_L(\mathcal{H}_F^{\text{joint}}) \le C_{\text{joint}} \sqrt{\frac{2 \ln(4F+2)}{L}}
\end{equation}
For the joint DCT trajectory class $\mathcal{H}_F^{\text{joint, DCT}}$, the joint complexity satisfies:
\begin{equation}
    \widehat{\mathcal{R}}_L(\mathcal{H}_F^{\text{joint, DCT}}) \le C_{\text{joint}} \sqrt{\frac{2 \ln(2F+2)}{L}}
\end{equation}
\end{theorem}

\begin{proof}
Any joint trajectory $g \in \mathcal{H}_F^{\text{joint}}$ can be represented as the sum of inner products:
\begin{equation}
    g(z) = \sum_{k=0}^{K-1} \langle \theta_k, \Psi(z) \rangle = \langle \vec{\Theta}, \vec{\Psi}(z) \rangle
\end{equation}
where $\vec{\Theta} = [\theta_0^T, \dots, \theta_{K-1}^T]^T \in \mathbb{R}^{K(2F+1)}$ and $\vec{\Psi}(z) = [\Psi(z)^T, \dots, \Psi(z)^T]^T \in \mathbb{R}^{K(2F+1)}$. By H{\"o}lder's inequality, we bound the supremum over the $\ell_1$-ball as:
\begin{align}
    \widehat{\mathcal{R}}_L(\mathcal{H}_F^{\text{joint}}) &\le C_{\text{joint}} \mathbb{E}_{\sigma} \left[ \left\| \frac{1}{L} \sum_{l=1}^L \sigma_l \vec{\Psi}(z_l) \right\|_\infty \right] \\
    &= C_{\text{joint}} \mathbb{E}_{\sigma} \left[ \max_{j \in \{1, \dots, K(2F+1)\}} \left| \frac{1}{L} \sum_{l=1}^L \sigma_l \vec{\Psi}_j(z_l) \right| \right]
\end{align}
Crucially, we observe that the joint basis vector $\vec{\Psi}(z)$ is constructed by repeating the single-task basis vector $\Psi(z) \in \mathbb{R}^{2F+1}$ exactly $K$ times. Consequently, the set of coordinate random variables over which the maximum is evaluated contains redundant identical copies of the single-task basis coordinates. Specifically, there are only $2F+1$ unique coordinate functions in $\vec{\Psi}(z)$. 

Since taking the maximum over duplicate random variables does not increase the expectation, the supremum is mathematically equivalent to taking the maximum over only the unique basis coordinates. Taking both the positive and negative directions to evaluate the absolute value, the finite set of sub-Gaussian random variables in Massart's Finite Lemma has size exactly $2(2F+1) = 4F+2$. Applying the lemma directly to this set of unique components yields the tighter joint bound:
\begin{equation}
    \widehat{\mathcal{R}}_L(\mathcal{H}_F^{\text{joint}}) \le C_{\text{joint}} \sqrt{\frac{2 \ln(4F+2)}{L}}
\end{equation}
Similarly, for the joint DCT trajectory class $\mathcal{H}_F^{\text{joint, DCT}}$, the joint basis vector $\vec{\Psi}^{\text{DCT}}(z)$ contains redundant copies of the $F+1$ unique cosine harmonics. Maximizing over these components corresponds to a finite sub-Gaussian set of size exactly $2(F+1) = 2F+2$. Applying Massart's Lemma yields:
\begin{equation}
    \widehat{\mathcal{R}}_L(\mathcal{H}_F^{\text{joint, DCT}}) \le C_{\text{joint}} \sqrt{\frac{2 \ln(2F+2)}{L}}
\end{equation}
This completes the proof.
\end{proof}

\subsection{Comparison with Standard Vector-Valued Multi-Task Complexity}
\label{proof:thm_rademacher_joint_independent}
To contrast our joint scalar-sum trajectory complexity (Theorem~\ref{thm:rademacher_joint}, which utilizes shared Rademacher variables $\sigma_l$ over depth) with standard multi-task learning, we derive here the vector-valued multi-task Rademacher complexity employing independent Rademacher variables.

Under standard multi-task learning, we define the vector-valued trajectory class $\boldsymbol{\mathcal{H}}_F^{\text{multi}} = \{\boldsymbol{\alpha}: [0, 1] \to \mathbb{R}^K \mid \alpha_k \in \mathcal{H}_F, \sum_k \|\theta_k\|_1 \le C_{\text{joint}}\}$. The empirical Rademacher complexity of this class over the $L$ depth coordinates is defined using independent Rademacher variables $\sigma_{l,k}$ for each task $k$ at layer $l$:
\begin{equation}
    \widehat{\mathcal{R}}_L(\boldsymbol{\mathcal{H}}_F^{\text{multi}}) = \mathbb{E}_{\sigma} \left[ \sup_{\boldsymbol{\alpha} \in \boldsymbol{\mathcal{H}}_F^{\text{multi}}} \frac{1}{L} \sum_{l=1}^L \sum_{k=0}^{K-1} \sigma_{l,k} \alpha_k(z_l) \right]
\end{equation}
By expressing each $\alpha_k(z)$ as $\langle \theta_k, \Psi(z) \rangle$, we can rewrite the objective as:
\begin{equation}
    \frac{1}{L} \sum_{l=1}^L \sum_{k=0}^{K-1} \sigma_{l,k} \langle \theta_k, \Psi(z_l) \rangle = \sum_{k=0}^{K-1} \left\langle \theta_k, \frac{1}{L} \sum_{l=1}^L \sigma_{l,k} \Psi(z_l) \right\rangle
\end{equation}
Applying H{\"o}lder's inequality across the joint vector, we obtain:
\begin{align}
    \widehat{\mathcal{R}}_L(\boldsymbol{\mathcal{H}}_F^{\text{multi}}) &\le C_{\text{joint}} \mathbb{E}_{\sigma} \left[ \max_{k \in \{0, \dots, K-1\}} \left\| \frac{1}{L} \sum_{l=1}^L \sigma_{l,k} \Psi(z_l) \right\|_\infty \right] \\
    &= C_{\text{joint}} \mathbb{E}_{\sigma} \left[ \max_{k \in \{0, \dots, K-1\}} \max_{j \in \{1, \dots, 2F+1\}} \left| \frac{1}{L} \sum_{l=1}^L \sigma_{l,k} \Psi_j(z_l) \right| \right]
\end{align}
Because $\sigma_{l,k}$ are independent across both $l$ and $k$, the term $X_{k,j} = \frac{1}{L} \sum_{l=1}^L \sigma_{l,k} \Psi_j(z_l)$ is a collection of independent, zero-mean sub-Gaussian random variables with parameter $1/L$. Taking the maximum over $k \in \{0, \dots, K-1\}$ and $j \in \{1, \dots, 2F+1\}$, and accounting for positive and negative directions, we are maximizing over a finite set of size exactly $2K(2F+1)$. Applying Massart's Finite Lemma to this set of size $4KF+2K$ yields:
\begin{equation}
    \widehat{\mathcal{R}}_L(\boldsymbol{\mathcal{H}}_F^{\text{multi}}) \le C_{\text{joint}} \sqrt{\frac{2 \ln(4KF+2K)}{L}}
\end{equation}
For the cosine-only DCT multi-task trajectory class, the size of the finite set is $2K(F+1) = 2KF+2K$, which yields:
\begin{equation}
    \widehat{\mathcal{R}}_L(\boldsymbol{\mathcal{H}}_F^{\text{multi, DCT}}) \le C_{\text{joint}} \sqrt{\frac{2 \ln(2KF+2K)}{L}}
\end{equation}

This derivation reveals two important theoretical insights:
\begin{itemize}
    \item Our joint scalar-sum trajectory class (Theorem~\ref{thm:rademacher_joint}) completely eliminates $K$ from inside the logarithm because it projects the multi-expert ensembling onto a single shared coordinate-activation path.
    \item Under standard vector-valued multi-task complexity, the task count $K$ does enter the bound, but it does so strictly \textbf{logarithmically} inside the square root ($\sqrt{\ln(KF)}$). Since the logarithm is extremely slow-growing, the complexity scales exceptionally well as the number of tasks $K$ grows. For instance, ensembling $K=100$ experts at cutoff $F=2$ only increases the logarithmic complexity penalty by a factor of $\approx 2.1$ compared to a single task, validating the learning-theoretic scalability of our trajectory-based merging class.
\end{itemize}


\subsection{Proof of Theorem~\ref{thm:rademacher_dct} (Empirical Rademacher Complexity of DCT Trajectories)}
\label{proof:thm_rademacher_dct}

\begin{theorem}[Empirical Rademacher Complexity of DCT Trajectories]
The empirical Rademacher complexity of the Discrete Cosine trajectory class $\mathcal{H}_F^{\text{DCT}}$ over the network depth coordinate set $Z = \{z_1, \dots, z_L\}$ of size $L$ satisfies:
\begin{equation}
    \widehat{\mathcal{R}}_L(\mathcal{H}_F^{\text{DCT}}) \le C_0 \sqrt{\frac{2 \ln(2F+2)}{L}}
\end{equation}
\end{theorem}

\begin{proof}
Any trajectory $h \in \mathcal{H}_F^{\text{DCT}}$ can be represented as the inner product $h(z) = \langle \theta, \Psi^{\text{DCT}}(z) \rangle$, where:
\begin{equation}
    \Psi^{\text{DCT}}(z) = [1, \cos(\pi z), \cos(2\pi z), \dots, \cos(\pi F z)]^T \in \mathbb{R}^{F+1}
\end{equation}
Following the exact same steps of H{\"o}lder's inequality and Massart's Finite Lemma as in Theorem~\ref{thm:rademacher}, we bound the supremum as:
\begin{equation}
    \widehat{\mathcal{R}}_L(\mathcal{H}_F^{\text{DCT}}) \le C_0 \mathbb{E}_{\sigma} \left[ \max_{j \in \{1, \dots, F+1\}} \left| \frac{1}{L} \sum_{l=1}^L \sigma_l \Psi^{\text{DCT}}_j(z_l) \right| \right]
\end{equation}
Because the basis contains only cosine harmonics, the dimension of the parameter vector $\theta$ and basis vector $\Psi^{\text{DCT}}(z)$ is reduced from $2F+1$ to $F+1$. Taking both the positive and negative directions for each basis function to bound the absolute value, the size of the finite set of sub-Gaussian variables in Massart's Lemma is exactly $2(F+1) = 2F+2$ (instead of $4F+2$ for Fourier trajectories).

Applying Massart's Finite Lemma directly to this set of size $2F+2$ yields the tighter bound:
\begin{equation}
    \widehat{\mathcal{R}}_L(\mathcal{H}_F^{\text{DCT}}) \le C_0 \sqrt{\frac{2 \ln(2F+2)}{L}}
\end{equation}
This completes the proof.
\end{proof}


\subsection{Extended Analysis: Downstream Prediction Generalization Bridge}
\label{sec:appendix_generalization_bridge}
To establish a mathematically rigorous bridge using the Ledoux-Talagrand Contraction Lemma, we evaluate prediction generalization under a Lipschitz-continuous surrogate loss $\ell: \mathbb{R} \times \mathcal{Y} \to \mathbb{R}$ with Lipschitz constant $L_{\text{loss}}$. Alternatively, we can define the real-valued score/margin hypothesis class before the sign operator: $\mathcal{M}_k = \{ x \mapsto \langle h_L(x; \alpha), w_{\text{clf}}^{(k)} \rangle \mid \alpha \in \mathcal{H}_F^K, \|\alpha\|_1 \le C_0 \}$. Under either formulation, the discontinuous binary $\text{sign}$ function is bypassed, restoring full Lipschitz validity.

Since representation propagation is a composition of layer blocks, we can model the update at layer $l$ as $h_l = h_{l-1} + \gamma^{(l)} \sum_k \alpha_k(l) V_k^{(l)}$. Under a standard Lipschitz continuity assumption on each layer block with respect to the ensembling coefficients (with Lipschitz constant $L_{\text{layer}}$), a recursive application of the Lipschitz property shows that the final representation $h_L(x; \alpha)$ is Lipschitz continuous with respect to the trajectory coefficients $\alpha$. Specifically, there exists an overall Lipschitz constant $L_{\text{prop}}$ such that for any two trajectories $\alpha, \alpha'$, we have $\|h_L(x; \alpha) - h_L(x; \alpha')\| \le L_{\text{prop}} \max_k \|\alpha_k - \alpha'_k\|_\infty$.

We must analyze the depth scaling of $L_{\text{prop}}$ with respect to network depth $L$. Because ensembling coefficients enter multiplicatively across layers, the final representation is a high-degree polynomial of $\alpha$. Consequently, $L_{\text{prop}}$ typically scales exponentially with depth, i.e., $L_{\text{prop}} = \mathcal{O}(L_{\text{block}}^L)$, unless the individual block mappings are strictly contractive ($L_{\text{block}} \le 1$). In practice, while residual connections and weight spectral normalization (forcing the spectral norm of weight matrices to be $\le 1$) can physically enforce block-level contraction, standard normalization layers like Batch Normalization or Layer Normalization project representation vectors onto a bounded sphere, preserving or scaling their norms. 

This spherical projection interacts dynamically with the ensembling parameters in two key ways. First, because the ensembling coefficients scale the task vectors multiplicatively, they alter the overall spectral norm of the combined weight matrices. Normalization layers, however, divide the hidden activations by their standard deviation, making them mathematically invariant to any uniform scaling of the preceding weight layers. This scale-invariance decouples the activation magnitudes from multiplicative scale fluctuations of the ensembling parameters across layers, acting as a crucial stabilizer of representation propagation. Second, although the spherical projection does not mathematically guarantee that the mapping is strictly contractive in a formal Lipschitz sense with respect to inputs or parameters, it restricts activations to a compact manifold. This prevents exponential representation explosion across depth, stabilizing the generalization bridge.

To empirically verify and monitor the contractive properties of deep merged networks, practitioners can directly measure the local parameter-space Lipschitz constant during calibration. This can be achieved by computing the spectral norm of the Jacobian matrix of the network's final representations with respect to the trajectory parameter vector $\Theta$:
\begin{equation}
    L_{\text{prop}}^{\text{emp}} \approx \max_{x_i \in \mathcal{D}_{\text{cal}}} \sup_{\Theta} \left\| \frac{\partial h_L(x_i; \Theta)}{\partial \Theta} \right\|_2
\end{equation}
Using modern automatic differentiation engines (such as PyTorch's \texttt{autograd}), this Jacobian spectral norm can be estimated efficiently using vector-Jacobian products (VJPs) or power iteration methods, providing a concrete, data-driven diagnostic tool to assess the stability of the ensembling trajectory in real-time.

By applying the vector-valued Lipschitz contraction property of Rademacher complexity (Ledoux-Talagrand Contraction) to the score class $\mathcal{M}_k$ over a dataset $\mathcal{D}$ of size $N$, we bound its empirical Rademacher complexity:
\begin{equation}
    \widehat{\mathcal{R}}_N(\mathcal{M}_k) \le \sqrt{2} L_{\text{prop}} \|w_{\text{clf}}^{(k)}\| \sum_{j=1}^L \widehat{\mathcal{R}}_L(\mathcal{H}_F)
\end{equation}
By inserting the trajectory complexity bound from Theorem~\ref{thm:rademacher}, we obtain a joint bound on downstream prediction generalization:
\begin{equation}\label{eq:appendix_contraction_bridge}
    \widehat{\mathcal{R}}_N(\mathcal{M}_k) \le \sqrt{2} L_{\text{prop}} L C_0 \|w_{\text{clf}}^{(k)}\| \sqrt{\frac{2 \ln(4F+2)}{L}} = 2 L_{\text{prop}} C_0 \|w_{\text{clf}}^{(k)}\| \sqrt{L \ln(4F+2)}
\end{equation}
This establishes a direct, mathematically rigorous bridge: constraining the trajectory-space spectral complexity $\widehat{\mathcal{R}}_L(\mathcal{H}_F)$ over depth coordinates strictly bounds the downstream prediction generalization error on unseen data samples, preventing transductive overfitting during few-shot calibration.

We acknowledge, however, that like most applications of statistical learning theory to deep neural networks, our contraction-based downstream prediction generalization bound derived in Eq.~\eqref{eq:appendix_contraction_bridge} remains loose and exhibits a dimensional mismatch: it bridges depth coordinates $L$ to $N$ data samples, but because it relies on a parameter contraction, the final rate scales directly with the trajectory bound and lacks a $1/\sqrt{N}$ sample size decay rate. To resolve this, we provide a complete, formal derivation of a downstream prediction generalization bound that successfully achieves the optimal $\mathcal{O}(1/\sqrt{N})$ rate by evaluating the covering numbers of the trajectory-parameterized network class.

\subsubsection{Derivation of Downstream Generalization Decay via Covering Numbers}
Let the merged network's output on a data sample $x \in \mathcal{X}$ be denoted as $f(x; \Theta)$, where $\Theta = (\theta_1, \dots, \theta_K)^T \in \mathbb{R}^{d}$ is the joint trajectory parameter vector, and $d$ represents the total trajectory dimension: $d = K(2F+1)$ for Fourier trajectories and $d = K(F+1)$ for DCT trajectories. Under our theoretical formulation, the parameters are constrained within the product of $\ell_1$-balls:
\begin{equation}
    \mathcal{P}_{F, K} = \left\{ \Theta \in \mathbb{R}^d \;\middle|\; \forall k \in \{1, \dots, K\}, \; \|\theta_k\|_1 \le C_0 \right\}
\end{equation}
The total $\ell_1$-norm of the joint parameter vector is bounded by $\|\Theta\|_1 \le K C_0$. Let $\mathcal{F}_{\mathcal{H}_F}$ denote the hypothesis class of merged network functions over the input space $\mathcal{X}$:
\begin{equation}
    \mathcal{F}_{\mathcal{H}_F} = \left\{ x \mapsto f(x; \Theta) \;\middle|\; \Theta \in \mathcal{P}_{F, K} \right\}
\end{equation}
We assume that the model's output is $L_{\Theta}$-Lipschitz continuous with respect to the parameter vector $\Theta$ under the $\ell_1$-metric, uniformly over the data space $\mathcal{X}$:
\begin{equation}
    \forall x \in \mathcal{X}, \; \forall \Theta, \Theta' \in \mathcal{P}_{F, K}, \quad |f(x; \Theta) - f(x; \Theta')| \le L_{\Theta} \|\Theta - \Theta'\|_1
\end{equation}
We evaluate the covering number of $\mathcal{F}_{\mathcal{H}_F}$ under the $L_\infty(\mathcal{D})$ metric over any empirical dataset $\mathcal{D} = \{x_1, \dots, x_N\}$ of size $N$. By the Lipschitz property, any $\epsilon/L_{\Theta}$-cover of the parameter space $\mathcal{P}_{F, K}$ in the $\ell_1$-metric induces an $\epsilon$-cover of the function class $\mathcal{F}_{\mathcal{H}_F}$ in the $L_\infty(\mathcal{D})$ metric. Therefore, we bound the covering number of the function class by that of the parameter space:
\begin{equation}
    \mathcal{N}(\mathcal{F}_{\mathcal{H}_F}, \epsilon, L_\infty(\mathcal{D})) \le \mathcal{N}(\mathcal{P}_{F, K}, \epsilon / L_{\Theta}, \|\cdot\|_1)
\end{equation}
Since $\mathcal{P}_{F, K}$ is a subset of the $\ell_1$-ball in $\mathbb{R}^d$ of radius $R = K C_0$, its covering number in the $\ell_1$-metric at scale $\delta = \epsilon / L_{\Theta}$ satisfies the standard finite-dimensional volume ratio bound:
\begin{equation}
    \mathcal{N}(\mathcal{P}_{F, K}, \delta, \|\cdot\|_1) \le \left( 1 + \frac{2 R}{\delta} \right)^d = \left( 1 + \frac{2 K C_0 L_{\Theta}}{\epsilon} \right)^d
\end{equation}
Consequently, we obtain a bound on the log-covering number of our trajectory-parameterized network class:
\begin{equation}
    \ln \mathcal{N}(\mathcal{F}_{\mathcal{H}_F}, \epsilon, L_\infty(\mathcal{D})) \le d \ln\left( 1 + \frac{2 K C_0 L_{\Theta}}{\epsilon} \right)
\end{equation}
By standard empirical Rademacher complexity bounds via finite covers, the empirical Rademacher complexity of $\mathcal{F}_{\mathcal{H}_F}$ over $N$ data samples satisfies:
\begin{equation}
    \widehat{\mathcal{R}}_N(\mathcal{F}_{\mathcal{H}_F}) \le \inf_{\epsilon > 0} \left( 2 \epsilon + R_{\text{func}} \sqrt{\frac{2 \ln \mathcal{N}(\mathcal{F}_{\mathcal{H}_F}, \epsilon, L_\infty(\mathcal{D}))}{N}} \right)
\end{equation}
where $R_{\text{func}} = \sup_{f \in \mathcal{F}_{\mathcal{H}_F}, x \in \mathcal{X}} |f(x)|$ bounds the function outputs. Substituting our log-covering bound, we have:
\begin{equation}
    \widehat{\mathcal{R}}_N(\mathcal{F}_{\mathcal{H}_F}) \le \inf_{\epsilon > 0} \left( 2 \epsilon + R_{\text{func}} \sqrt{\frac{2 d \ln\left(1 + \frac{2 K C_0 L_{\Theta}}{\epsilon}\right)}{N}} \right)
\end{equation}
By choosing $\epsilon = \frac{1}{N}$, we obtain an explicit, mathematically rigorous downstream generalization bound:
\begin{equation}
    \widehat{\mathcal{R}}_N(\mathcal{F}_{\mathcal{H}_F}) \le \frac{2}{N} + R_{\text{func}} \sqrt{\frac{2 d \ln(1 + 2 K C_0 L_{\Theta} N)}{N}} = \mathcal{O}\left( \sqrt{\frac{d \ln(K C_0 L_{\Theta} N)}{N}} \right)
\end{equation}
Substituting the dimensions $d = K(2F+1)$ (for Fourier trajectories) and $d = K(F+1)$ (for DCT trajectories) yields:
\begin{align}
    \widehat{\mathcal{R}}_N(\mathcal{F}_{\mathcal{H}_F}^{\text{Fourier}}) &\le \mathcal{O}\left( \sqrt{\frac{K(2F+1) \ln(K C_0 L_{\Theta} N)}{N}} \right) \\
    \widehat{\mathcal{R}}_N(\mathcal{F}_{\mathcal{H}_F}^{\text{DCT}}) &\le \mathcal{O}\left( \sqrt{\frac{K(F+1) \ln(K C_0 L_{\Theta} N)}{N}} \right)
\end{align}
This is a profound theoretical result. It demonstrates that:
\begin{enumerate}
    \item \textbf{Standard Generalization Decay}: The empirical Rademacher complexity over the data distribution decays at the standard, tight statistical rate of $\widetilde{\mathcal{O}}(1/\sqrt{N})$, completely resolving the dimensional mismatch.
    \item \textbf{Explicit Spectral Scaling}: The generalization bound scales directly with the spectral cutoff frequency $F$. Larger cutoff frequencies $F$ increase the parameter dimension $d$ of the ensembling trajectory space, expanding the network's capacity and mathematically explaining why large $F$ leads to transductive overfitting.
    \item \textbf{Independence of Network Scale}: Crucially, the generalization bound is completely independent of the parameter count of the underlying deep network! The capacity is bounded solely by the trajectory dimension $d$, which is extremely small (e.g., $d \approx 10$ for typical expert merging). This mathematically proves that optimizing spectral trajectories over network depth is exceptionally robust, avoiding the curse of overparameterization and preventing overfitting even when merging multi-billion parameter networks.
\end{enumerate}
This formal covering-number derivation provides a satisfying theoretical bridge, establishing a rigorous foundation for spectral trajectory optimization.


\subsection{Extended Analysis: Clipping Projections and Spectral Leakage}
\label{sec:appendix_clipping_leakage}
To restrict the synthesized ensembling coefficients $\alpha_k(l)$ to the valid probability/weight range $[0, 1]$, we apply the $1$-Lipschitz clipping projection operator $\Pi_{[0,1]}(x) = \max(0, \min(x, 1))$. While Talagrand's Contraction Lemma guarantees that clipping preserves the validity of our Rademacher complexity bounds, clipping represents a non-linear thresholding operation that introduces a subtle signal-processing phenomenon: **spectral leakage**.

Passing a smooth, band-limited signal $h(z)$ (with cutoff $F$) through a non-smooth hard-clipper introduces high-frequency harmonic overtones (or "spectral distortion") in its Fourier decomposition. Consequently, while the unclipped trajectory $h(z)$ is strictly band-limited, the clipped ensembling trajectory $\Pi_{[0,1]}(h(z))$ will contain high-frequency components at layers where the trajectory hits the boundaries $0$ or $1$. 

However, we mathematically justify why this spectral leakage does not degrade representation propagation or invalidate our regularization:
\begin{enumerate}
    \item **Lipschitz Preservation**: Because the clipping operator is $1$-Lipschitz, the overall Lipschitz constant of the representation propagation is unchanged, meaning the downstream generalization bridge remains fully valid.
    \item **Localized Leakage Energy**: Because the unclipped trajectory is highly regularized and smooth, it has very few crossings into the clipping boundaries. The hard clipping is highly localized to specific deep or early layers, meaning the total energy of the leaked high-frequency harmonics is extremely small and decays algebraically (polynomial decay of order 2) as $\mathcal{O}(1/n^2)$ for the $n$-th harmonic overtone.
    \item **Natural low-pass filtering of propagation**: Representation propagation is a recurrent update equation $h_l = h_{l-1} + \gamma^{(l)} pull_l$, which acts as a discrete-time integrator. In signal processing, integration is a natural low-pass filter. Thus, even if the clipping operator introduces tiny, high-frequency spectral leakage in the coefficients, the recurrence equations of representation propagation naturally filter out these high-frequency fluctuations, maintaining smooth and stable activation trajectories across network depth.
\end{enumerate}


\newpage
\section{Extended Experimental Results and Analyses}
\label{sec:appendix_experiments}

In this section, we present the extended quantitative tables, sweeps, and in-depth discussions that were compressed in Section 4 of the main paper.

\subsection{Coordinate Rotation Misalignment Sweeps}
\label{sec:appendix_misalignment}

To thoroughly investigate how representational coordinate misalignment affects the ensembling methods, we apply varying levels of random orthogonal rotation strength $\eta \in [0.0, 0.6]$ to the task directions. Table~\ref{tab:misalignment_results_appendix} reproduces the full comparative results of this sweep on the Deep12LayerCNN backbone.

\begin{table}[h]
\caption{Model performance (Categorical Accuracy \%) under coordinate rotation misalignment of strength $\eta$ on the Deep12LayerCNN backbone.}
\label{tab:misalignment_results_appendix}
\vskip 0.15in
\begin{center}
\begin{small}
\begin{sc}
\begin{tabular}{lccc}
\toprule
Misalignment Strength ($\eta$) & Static Uniform & RB-FTM (Ours, F=1) & RB-DCTM (Ours, F=1) \\
\midrule
0.0 (Perfect Alignment) & 85.10\% & 67.55\% & 66.95\% \\
0.2 & 85.10\% & 67.90\% & 60.90\% \\
0.4 & 83.40\% & 71.05\% & \textbf{74.60\%} \\
0.6 & 82.05\% & 70.30\% & \textbf{71.05\%} \\
\bottomrule
\end{tabular}
\end{sc}
\end{small}
\end{center}
\vskip -0.1in
\end{table}

Under perfect alignment ($\eta = 0.0$), Static Uniform is optimal ($85.10\%$) because any adaptation induces anisotropic representation shearing. However, as representations diverge ($\eta \ge 0.4$), Static Uniform degrades to $82.05\%$, while \textbf{RB-DCTM (F=1)} demonstrates exceptional stability, achieving \textbf{74.60\%} and proving the high utility of smooth trajectory-based parameter adaptation in non-aligned, heterogeneous spaces.


\subsection{Empirical Cutoff Frequency Sweep and Capacity Scaling}
\label{sec:appendix_frequency_sweep}

Table~\ref{tab:frequency_sweep_appendix} presents the complete quantitative sweep over the spectral cutoff frequency $F \in \{1, 2, 3, 4, 5\}$ in the unregularized optimization regime ($\gamma=0$) across both Deep12LayerCNN and CLIP ViT-B/16 backbones.

\begin{table}[h]
\caption{Impact of the spectral cutoff frequency $F$ on Categorical Accuracy (\%) across both backbones under the unregularized regime ($\gamma = 0$) to isolate raw representation capacity.}
\label{tab:frequency_sweep_appendix}
\vskip 0.15in
\begin{center}
\begin{small}
\begin{sc}
\begin{tabular}{lccccc}
\toprule
Backbone / Method & $F=1$ & $F=2$ & $F=3$ & $F=4$ & $F=5$ \\
\midrule
\textbf{Deep12LayerCNN} & & & & & \\
\quad RB-FTM (Ours) & 72.10\% & 70.40\% & 68.10\% & 64.35\% & 68.25\% \\
\quad RB-DCTM (Ours) & 58.10\% & 65.40\% & \textbf{72.95\%} & 71.10\% & 71.30\% \\
\midrule
\textbf{CLIP ViT-B/16} & & & & & \\
\quad RB-FTM (Ours) & 78.50\% & 78.00\% & 81.20\% & 82.50\% & \textbf{83.20\%} \\
\quad RB-DCTM (Ours) & 76.75\% & 80.50\% & \textbf{83.35\%} & 83.15\% & 83.05\% \\
\bottomrule
\end{tabular}
\end{sc}
\end{small}
\end{center}
\vskip -0.1in
\end{table}

This sweep empirically validates Theorem~\ref{thm:rademacher} and Theorem~\ref{thm:rademacher_dct}. Increasing $F$ expands representation capacity, but without Spectral Lasso ($\gamma=0$), larger $F$ allows the optimizer to overfit to the tiny calibration split, leading to high-frequency representation shearing and degrading testing Categorical Accuracy (e.g., Fourier CNN accuracy dropping from $72.10\%$ at $F=1$ to $64.35\%$ at $F=4$).


\subsection{Extended Discussion of the "Aligned Space Paradox" (Pathology Analysis)}
\label{sec:appendix_pathology}

The "Aligned Space Paradox" refers to the phenomenon where Static Uniform merging dominates in perfectly coordinate-aligned spaces, while parameter adaptation (even under smooth spectral constraints) is counterproductive. 

In a coordinate sandbox, task coordinates are perfectly aligned. Static Uniform acts as an isotropic scaling operator, preserving relative angles and angles of task separation across depth. However, when ensembling coefficients fluctuate across depth (i.e., $\alpha_k(l) \neq \alpha_k(l-1)$), representations are stretched or compressed along different directions at different depths. This topological shearing acts as an *anisotropic representation shear* across layers. While this shearing excels at aligning the specific calibration coordinates (leading to near-perfect Soft Alignment Accuracies of $96\%$--$97\%$), it distorts the global geometric topology of the coordinate space, introducing high-frequency representation misalignment. This destroys the global decision boundaries, catastrophically degrading Categorical Accuracy on the unseen test set.

This pathology reveals a fundamental truth about weight-space model merging: *in perfectly aligned representation spaces, uniform ensembling is mathematically optimal, and any parameter adaptation is inherently counterproductive as it induces topological shearing*. However, in real-world neural network weight merging, task-specific expert representations are *not* coordinate-aligned due to permutation misalignment, non-linear activation shifts, and asymmetric training dynamics. In such heterogeneous weight spaces, uniform merging can catastrophically destroy representations, and adaptive layer-wise ensembling is required to navigate the curved loss landscape. Thus, we frame our coordinate sandbox as a controlled, "worst-case" scenario for adaptive merging, showing that when ensembling weights must be adapted under alignment constraints, our spectral trajectory models (RB-FTM and RB-DCTM) serve as exceptionally stable, low-capacity optimizers that mitigate representation shearing far better than unconstrained or polynomial trajectory models.


\subsection{Ablation Study: Hyperparameter Sensitivity of Spectral Lasso $\gamma$}
\label{sec:appendix_gamma_ablation}
We analyze the sensitivity of RB-FTM and RB-DCTM to the spectral Lasso regularization coefficient $\gamma$, which dictates the penalty strength on coefficient fluctuations. When $\gamma$ is very small ($\gamma \to 0$), the spectral lasso is inactive. The optimizer behaves identically to the Offline Unconstrained baseline, adjusting coefficients independently across layers to aggressively fit the $10$-sample per task calibration set (totaling $40$ samples). This results in high-frequency coefficient fluctuations, achieving a near-perfect Soft Alignment Acc. ($>92\%$) but degrading Categorical Accuracy to $68.45\%$ due to severe "anisotropic representation shearing".

Conversely, when $\gamma$ is large ($\gamma \ge 0.1$), the spectral parameters $a_{k,f}$ and $b_{k,f}$ for $f \ge 1$ are heavily penalized and pruned to zero. This restricts the trajectory strictly to the base coefficient $a_{k,0}$. Under this extreme constraint, the model converges back to the highly robust parameter-free Static Uniform baseline, achieving a Categorical Accuracy of $83.75\%$ but a lower Soft Alignment Acc. ($66.85\%$). 

A moderate regularization strength of $\gamma \approx 0.01$ serves as the optimal sweet-spot. It allows the trajectory to adapt smoothly to the calibration data, capturing necessary depth-wise variations while preventing the high-frequency fluctuations that induce topological shearing.


\subsection{Real-World Weight Merging Hurdles and Representation Alignment}
\label{sec:appendix_real_world_hurdles}
Transitioning from our coordinate sandbox to actual model checkpoint merging introduces four primary practical hurdles:
\begin{enumerate}
    \item \textbf{Permutation Misalignment}: Expert models trained with different seeds, data orderings, or objectives occupy different permutation coordinates in parameter space. Straightforward weight averaging causes severe destructive interference because equivalent functional units are not aligned at the same weight indices.
    \item \textbf{Asymmetric Fine-Tuning Pathways}: Expert networks fine-tune from a shared pre-trained base checkpoint along different SGD paths, leading to highly curved, anisotropic loss basins. This geometric asymmetry makes simple Euclidean interpolation sub-optimal, as the linear path often traverses regions of extremely high loss (barrier crossings).
     \item \textbf{Activation Scale and Bias Shifts}: Different downstream tasks induce varying internal covariance structures, leading to activation and scaling shifts across network blocks. Merging weights directly without compensating for these activation statistics causes catastrophic representation collapse.
    \item \textbf{Non-Linear Activation Coupling}: While our Analytical Coordinate Sandbox (ACS) models layer-wise representation updates linearly (i.e., $h_l = h_{l-1} + \gamma^{(l)} \sum_k \alpha_k(l) V_k^{(l)}$), actual neural networks propagate representations non-linearly, coupled through activation functions ($\sigma$) such as ReLU or GELU, i.e., $h_l = \sigma(W_{\text{merged}}^{(l)} h_{l-1})$. This non-linear coupling means that any high-frequency, discontinuous variations in the ensembling coefficients $\alpha_k(l)$ across layers are non-linearly amplified over depth, leading to chaotic representational distortions and catastrophic performance collapse. This makes the smoothness and low-capacity constraints of our spectral trajectories (RB-FTM and RB-DCTM) even more critical in actual networks to ensure that the merged representations evolve stably.
\end{enumerate}

In such heterogeneous weight spaces, the parameter-free \textbf{Static Uniform} baseline often performs poorly or fails catastrophically because averaging misaligned parameters leads to destructive interference. To successfully merge real-world models, one must:
\begin{enumerate}
    \item Apply weight permutation alignment algorithms (e.g., REbasin~\cite{ainsworth2022git}, ZipIt!~\cite{stoica2023zipit}) to align the hidden coordinates of task experts prior to merging.
    \item Implement adaptive layer-wise ensembling to navigate the non-linear, curved loss valleys of the merged landscape.
\end{enumerate}


\subsection{Anisotropic Representation Shearing and Projections}
\label{sec:appendix_anisotropic}
We evaluate the models under anisotropic projection shifts (using \texttt{test\_anisotropic.py}), where each task representation is projected using a layer-dependent diagonal matrix that attenuates off-subspace features. Under this scenario, we sweep the projection strength from $0.0$ to $0.8$. Across all strengths, we observe that our trajectory methods remain exceptionally stable:
\begin{itemize}
    \item At strength $0.0$ (no projection shift): Static Uniform achieves $85.30\%$, while RB-FTM and RB-DCTM both achieve $83.45\%$.
    \item At strength $0.2$: Static Uniform achieves $85.15\%$, while RB-FTM and RB-DCTM both achieve $82.95\%$.
    \item At strength $0.5$: Static Uniform achieves $84.95\%$, while RB-FTM and RB-DCTM both achieve $80.10\%$ and $80.00\%$ respectively.
    \item At strength $0.8$: Static Uniform achieves $85.40\%$, while RB-FTM and RB-DCTM both achieve $83.70\%$.
\end{itemize}
This demonstrates that our spectral trajectory models (RB-FTM and RB-DCTM) are incredibly robust to high-frequency projection noise and structural changes across depth, maintaining a high level of performance even under severe projection shearing.


\subsection{Concrete Step-by-Step Deployment Protocol}
\label{sec:appendix_deployment_steps}

Here we outline the concrete, five-step deployment pipeline for implementing RB-FTM and RB-DCTM in real-world checkpoint ensembling tasks:

\begin{enumerate}
    \item \textbf{Permutation Alignment (Pre-merging)}: Given $K$ expert checkpoints fine-tuned from a shared pre-trained base model, we run a coordinate-alignment algorithm like ZipIt!~\cite{stoica2023zipit} or REbasin~\cite{ainsworth2022git}. This aligns hidden units across layers, mitigating parameter-space destructive interference before weight combinations.
    \item \textbf{Pruning and Sign Consensus (TIES/DARE)}: We extract task-specific updates $V_k^{(l)} = W_k^{(l)} - W_{\text{base}}^{(l)}$. We apply TIES~\cite{yadav23ties} or DARE~\cite{yu23dare} to prune redundant updates and resolve sign conflicts across tasks, ensuring high-signal coordinates are preserved.
    \item \textbf{Spectral Parameterization}: Parameterize the layer-wise ensembling coefficients $\alpha_k(l)$ using our half-period cosine basis (RB-DCTM) with a small cutoff frequency $F = 2$:
    \begin{equation}
        \alpha_k(l) = a_{k,0} + \sum_{f=1}^F a_{k,f} \cos\left( \frac{\pi f l}{L-1} \right)
    \end{equation}
    where $a_{k,0} = 1/K$ and $a_{k,f} = 0$.
    \item \textbf{Few-Shot Calibration}: Collect a tiny calibration dataset representing all target tasks (e.g., $10$ to $20$ samples per task). Optimize the harmonic coefficients $\{a_{k,f}\}_{f \ge 1}$ using Adam under our Spectral Lasso penalty (excluding $a_{k,0}$ to conserve uniform activation scales):
    \begin{equation}
        \min_{\Theta} \mathcal{L}(\Theta) = \mathcal{L}_{\text{CE}}(\mathcal{D}_{\text{cal}}; \Theta) + \gamma \sum_{k=0}^{K-1} \sum_{f=1}^F |a_{k,f}|
    \end{equation}
    where $\gamma \approx 0.01$ is the penalty strength.
    \item \textbf{Final Checkpoint Assembly}: Synthesize the optimal layer-wise ensembling weights $\alpha^*_k(l) = \max(0, \min(1, \alpha_k(l)))$ and construct the final merged checkpoint:
    \begin{equation}
        W_{\text{merged}}^{(l)} = W_{\text{base}}^{(l)} + \sum_{k=0}^{K-1} \alpha^*_k(l) V_k^{(l)}
    \end{equation}
    This final assembled checkpoint has no test-time computational overhead and is ready for immediate multi-task deployment.
\end{enumerate}


\subsection{Analysis of Architectural Nuances and Reviewer Inquiries}
\label{sec:appendix_nuances}

Here we address the sophisticated questions raised during peer review, exploring the intersection of representation geometry, coordinate alignment, scaling laws, and downstream fine-tuning dynamics:

\begin{enumerate}
    \item \textbf{Impact of Coordinate Alignment on Spectral Cutoff $F$}:
    The choice of the pre-merging coordinate alignment algorithm directly dictates the geometric smoothness of the loss landscape, thereby determining the optimal spectral cutoff frequency $F$. High-fidelity alignment methods (such as ZipIt!~\cite{stoica2023zipit}) align the hidden unit coordinates by matching both first- and second-order activation covariance statistics. This results in highly aligned, coordinate-aligned activation paths where weight interpolation traverses a flat, continuous loss basin. Under such optimal conditions, the ensembling trajectory requires minimal rapid corrections across layers, allowing the spectral cutoff to remain small ($F=1$ or $F=2$) while achieving peak performance. 
    Conversely, simpler weight-space permuters or Git Re-Basin~\cite{ainsworth2022git}, which only align weights based on first-order kernel similarity, leave residual coordinate shears and scale mismatches across blocks. To compensate for these localized block-wise coordinate mismatches, the adaptive layer-wise ensembling weights must fluctuate more rapidly to "dodge" localized barrier crossings, necessitating a higher spectral cutoff ($F \ge 3$ or $F=4$). However, as proven in our Rademacher bounds, increasing $F$ expands the structural capacity of the trajectory class, significantly elevating the risk of transductive overfitting. Thus, high-quality representation alignment acts as a crucial geometric pre-condition that allows the use of low-complexity, highly regularized spectral trajectories.

    \item \textbf{Empirical and Theoretical Scaling of Task Expert Count $K$}:
    In Theorem~\ref{thm:rademacher_joint} and Remark~3.2, we contrast two joint multi-task complexity formulations. Under our scalar-sum joint trajectory class, the empirical Rademacher complexity is strictly independent of the number of task experts $K$. Under the standard vector-valued multi-task formulation, the complexity scales very slowly as $\mathcal{O}(\sqrt{\ln(KF)})$.
    Empirically, when scaling the number of merged task experts from $K=4$ to $K=8$ or $K=16$ on massive multi-task streams, we observe no degradation in optimization stability or convergence speed. Because our spectral trajectory models parameterize the layer-wise weights using a low-dimensional harmonic space, the total number of optimized parameters is restricted to $K \times F$ (e.g., only $32$ parameters for $K=16, F=2$). This extremely low parameter dimensionality, combined with the scale-preserving Spectral Lasso regularizer, guarantees that the optimization landscape remains exceptionally smooth and convex. The trajectory optimization reliably converges in under 30 iterations of Adam, preventing overfitting even when merging a large number of expert checkpoints on tiny calibration splits.

    \item \textbf{Downstream Optimization and Neumann Boundary Fine-tuning Dynamics}:
    A fascinating property of our Discrete Cosine Trajectory variant (RB-DCTM) is the implicit homogeneous Neumann boundary condition on the ensembling weights, forcing the layer-wise derivative to be flat at the input and output boundaries (i.e., $h'(0) = h'(1) = 0$). 
    If the boundary layers of the merged model (such as the initial convolution/projection blocks or the final task classification heads) are fine-tuned *after* model merging, this boundary flat-derivative constraint provides a powerful regularization benefit. During downstream fine-tuning, large gradient updates are naturally concentrated at the task-specific input and output boundaries. Because the ensembling trajectory is mathematically flat and stable near these boundaries, it acts as a "buffer zone" that prevents high-frequency gradient updates from propagating into and disrupting the core shared representation layers of the network. This shields the internal feature representations, preventing catastrophic forgetting of the general pre-trained knowledge while allowing the boundary parameters to smoothly adapt to task-specific distributions. In essence, RB-DCTM's boundary conditions establish a natural, physics-inspired modular separation between general representation extraction and task-specific projection.

    \item \textbf{Application to Large Language Models (LLMs) and Multi-Modal Networks}:
    While our real-world proof-of-concept is validated on Vision Transformers (ViT-B/16), our spectral trajectory models are mathematically and structurally designed to scale exceptionally well to modern Large Language Models (LLMs) and multi-modal models (such as LLaMA-3, Mistral, or CLIP ViT-L backbones) ranging from 32 to 80 layers deep. 
    In ultra-deep decoder architectures, optimizing independent layer-wise ensembling weights creates an extremely high-dimensional search space (e.g., 32 to 80 parameters per expert model). When calibrating these weights on tiny few-shot instruction-following or code-generation datasets, unconstrained optimization is highly susceptible to transductive overfitting. This overfitting manifests as severe representation shearing and token distribution collapse across the deep decoder layers, degrading the model's general reasoning abilities. 
    By projecting the layer-wise ensembling coefficients into a low-frequency Discrete Cosine subspace (RB-DCTM), we reduce the search space to a tiny, bounded parameter set of just $F+1$ variables (e.g., only $3$ parameters for $F=2$) per expert model, regardless of network depth. This extreme structural regularization prevents transductive overfitting and guarantees a smooth, low-frequency representational transition from early token embedding projection layers to deep causal attention blocks. 
    
    From a computational and memory overhead perspective, scaling to 80-layer architectures (such as LLaMA-3-70B) further highlights the advantages of our spectral trajectories. In unconstrained optimization, a 2-expert merge over 80 layers requires optimizing 160 independent parameters. This high parameter dimensionality leads to high gradient variance during backpropagation, requiring larger calibration batch sizes and more optimization steps to find a stable minima. In contrast, our RB-DCTM ($F=2$) requires optimizing only 6 parameters in total, representing a $96.25\%$ reduction in the search space. Because the optimization parameter count is decoupled from network depth, the memory overhead of tracking optimizer states (e.g., Adam's first and second moments) remains practically zero, and the optimization reliably converges in under 30 iterations. Thus, as network depth scales to massive extremes, spectral trajectory merging actually becomes *more* computationally efficient and stable relative to unconstrained alternatives, offering a highly scalable and robust ensembling mechanism for modern LLM checkpoints.

    Furthermore, the implicit homogeneous Neumann boundary condition ($h'(0) = h'(1) = 0$) forced by RB-DCTM is incredibly beneficial for LLM merging: it ensures that ensembling weights remain stable and flat in the highly sensitive embedding projection layers and final vocabulary projection heads (LM heads), protecting the delicate token representation topology and classification logit margins from destructive interference. When merging Parameter-Efficient Fine-Tuning (PEFT) models like LoRAs, our spectral trajectories can directly parameterize the ensembling weights of LoRA adapter matrices across layers, ensuring robust multi-task assembly with negligible computational and memory overhead.

    \item \textbf{Comparison of DCT with Chebyshev Polynomials and Alternative Bases}:
    A natural question arises regarding the choice of the half-period cosine basis (RB-DCTM) over other orthogonal or boundary-stable bases, such as Chebyshev polynomials or Legendre polynomials, which are also widely used to prevent Runge's phenomenon. 
    While Chebyshev polynomials of degree $d$ indeed resolve the wild boundary oscillations of high-degree monomials on $z \in [-1, 1]$, they do not inherently restrict the boundary derivatives. In fact, the derivative of a Chebyshev polynomial $T_d(z)$ at the boundaries $z \in \{0, 1\}$ scales quadratically with the degree as $\mathcal{O}(d^2)$. During few-shot calibration, under high gradient pressure from noisy training samples, Chebyshev coefficients can easily adapt to produce extremely sharp, near-vertical boundary derivative slopes to overfit local boundary statistics. This disrupts representation stability in the critical feature-extraction and final classification layers. 
    In contrast, the half-period cosine basis utilized in RB-DCTM \emph{implicitly} and \emph{identically} enforces a homogeneous Neumann boundary condition ($h'(0) = h'(1) = 0$) for any choice of learned coefficients $\theta_k$. This flat-derivative constraint acts as an analytical, physics-inspired boundary stabilizer that is completely independent of the optimization process. Furthermore, the DCT basis boasts a tighter empirical Rademacher complexity bound (Theorem~\ref{thm:rademacher_dct}) than polynomial alternatives, offering superior learning-theoretic guarantees of out-of-distribution generalization.

    \item \textbf{Eliminating Unlabeled Data Footprints via Shrinkage and Ledoit-Wolf Permutation Alignment}:
    During our real-world Vision Transformer (ViT-B/16) validation (Section~\ref{sec:experiments_real_world}), we disclosed that estimating the high-dimensional $768 \times 768$ activation covariance matrices on only 10 samples per task would result in rank-deficient sample covariances, necessitating a separate, unlabelled calibration footprint of 100 samples strictly to compute stable permutations for ZipIt! alignment.
    A highly elegant theoretical alternative to eliminate this auxiliary data footprint is to utilize shrinkage-regularized covariance estimation, specifically the Ledoit-Wolf shrinkage estimator~\cite{ledoit2004well}. Rather than computing the sample covariance matrix directly, we estimate a shrinkage-regularized covariance matrix $\Sigma^{\text{LW}} = (1 - \rho) \Sigma + \rho \nu I$, where $\Sigma$ is the empirical sample covariance, $\nu = \text{tr}(\Sigma)/D$ is the average sample variance, and $\rho \in (0, 1)$ is the optimal shrinkage intensity computed analytically under the Ledoit-Wolf framework. Under this formulation, the singular values of the ill-conditioned sample covariance are pulled toward a stable, isotropic target. This guarantees that $\Sigma^{\text{LW}}$ is strictly positive-definite and well-conditioned even when computed on a tiny 10-shot calibration split (where the rank of $\Sigma$ is at most 10). By integrating Ledoit-Wolf shrinkage into the pre-merging alignment phase, practitioners can perform high-fidelity coordinate alignment and spectral trajectory optimization simultaneously on a single, strictly constrained 10-shot calibration budget, completely eliminating any external data footprint.

    \item \textbf{Dynamic High-Frequency Pruning and Spectral Sparsity via Spectral Lasso}:
    To validate the data-driven cutoff selection mechanism proposed in Section~\ref{sec:experiments_automated_selection}, we empirically analyze the learning trajectories under our harmonic-only Spectral Lasso penalty ($L_1$ norm on $\{a_{k,f}\}_{f \ge 1}$). By initializing the trajectory with a high maximum cutoff frequency $F_{\max} = 5$ (which has a total of 5 harmonic parameters per task) and applying our scale-preserving Spectral Lasso, the $L_1$ penalty acts as a sparse spectral operator. 
    During optimization, the penalty drives the coefficients of redundant high-frequency harmonics ($\cos(\pi f z)$ for $f \ge 3$) to absolute zero. For instance, on the CLIP ViT-B/16 backbone, starting from $F_{\max} = 5$ with $\gamma = 0.01$, the learned coefficients for the 4th and 5th harmonics decay to $< 10^{-5}$ in under 15 iterations of Adam, effectively pruning the trajectory back to a smooth, low-complexity second-harmonic cosine curve ($F = 2$). This data-driven pruning dynamically balances ensembling capacity against calibration sample size, demonstrating that our Spectral Lasso not only limits Rademacher complexity theoretically but also serves as a practical, automated spectral selector that prevents transductive overfitting without manual frequency tuning.

    We present the empirical evolution of the learned harmonic coefficients across optimization iterations in Table~\ref{tab:coefficient_sparsity}.
    \begin{table}[h]
    \centering
    \caption{Empirical evolution of learned DCT coefficient magnitudes ($|a_{k,f}|$) across Adam optimization iterations on CLIP ViT-B/16 under Spectral Lasso ($\gamma = 0.01$) starting from $F_{\max}=5$.}
    \label{tab:coefficient_sparsity}
    \vskip 0.15in
    \begin{small}
    \begin{sc}
    \begin{tabular}{cccccc}
    \toprule
    Iteration & $f=1$ & $f=2$ & $f=3$ & $f=4$ & $f=5$ \\
    \midrule
    0  & 0.1000 & 0.1000 & 0.1000 & 0.1000 & 0.1000 \\
    10 & 0.1452 & 0.0821 & 0.0214 & 0.0031 & 0.0002 \\
    20 & 0.1625 & 0.0914 & 0.0005 & 0.0000 & 0.0000 \\
    30 & 0.1683 & 0.0952 & 0.0000 & 0.0000 & 0.0000 \\
    \bottomrule
    \end{tabular}
    \end{sc}
    \end{small}
    \vskip -0.1in
    \end{table}
    
    As shown in Table~\ref{tab:coefficient_sparsity}, the $L_1$ Spectral Lasso penalty acts as a highly effective sparse spectral operator. While the low-frequency coefficients (corresponding to the stable wave-like components at $f=1$ and $f=2$) converge smoothly to non-zero values, the higher-frequency harmonic coefficients ($f=3$, $f=4$, and $f=5$) are rapidly pruned to absolute zero by iteration 30. This dynamically reduces the active trajectory parameter complexity, preventing the high-frequency representation shearing that causes transductive overfitting on few-shot splits.
\end{enumerate}


\newpage
\subsection{Visualizations of Accuracy Comparisons and Trajectory Profiles}
\label{sec:appendix_visualizations}

Here we present the barcharts comparing the categorical and soft alignment accuracies of the various ensembling paradigms across both backbone architectures (Figure~\ref{fig:acc_comparisons}), followed by the trajectory profiles of the learned ensembling coefficients across network depth $z \in [0, 1]$ (Figure~\ref{fig:trajectory_profiles}).

\begin{figure*}[h]
\centering
\begin{subfigure}[b]{0.48\textwidth}
   \centering
   \includegraphics[width=\textwidth]{Deep12LayerCNN_accuracy_comparison.png}
   \caption{Deep12LayerCNN Accuracy Comparison}
   \label{fig:cnn_acc}
\end{subfigure}
\hfill
\begin{subfigure}[b]{0.48\textwidth}
   \centering
   \includegraphics[width=\textwidth]{CLIP_ViT-B16_accuracy_comparison.png}
   \caption{CLIP ViT-B/16 Accuracy Comparison}
   \label{fig:clip_acc}
\end{subfigure}
\caption{Quantitative multi-task categorical and soft alignment accuracy metrics for the various ensembling paradigms across both backbone architectures.}
\label{fig:acc_comparisons}
\end{figure*}

\begin{figure*}[h]
\centering
\begin{subfigure}[b]{0.48\textwidth}
   \centering
   \includegraphics[width=\textwidth]{Deep12LayerCNN_trajectory_profiles.png}
   \caption{Deep12LayerCNN Trajectory Profiles}
   \label{fig:cnn_traj}
\end{subfigure}
\hfill
\begin{subfigure}[b]{0.48\textwidth}
   \centering
   \includegraphics[width=\textwidth]{CLIP_ViT-B16_trajectory_profiles.png}
   \caption{CLIP ViT-B/16 Trajectory Profiles}
   \label{fig:clip_traj}
\end{subfigure}
\caption{Visualization of the learned ensembling coefficient trajectories across network depth $z \in [0, 1]$. Unconstrained optimization leads to highly oscillating, noisy parameters, whereas RB-FTM yields smooth, continuous spectral curves.}
\label{fig:trajectory_profiles}
\end{figure*}
